% !PRINTED 592
\weberpage{592}

\begin{center}
\S.\,174.

\gesperrt{Der K\"orper der dritten Einheitswurzeln.}
\end{center}

Der Hauptsatz des vorigen Paragraphen, dass jede ganze Zahl des K\"orpers $R(i)$ sich nur auf eine Art in unzerlegbare Factoren zerlegen l\"asst, demzufolge der Begriff der unzerlegbaren Zahl mit dem der Primzahl zusammenf\"allt, beruht, wie man sieht, in der Hauptsache auf dem Algorithmus des gr\"ossten gemeinschaftlichen Theilers, der in den Formeln \S.\,173, (1) enthalten ist, und wenn man also in einem anderen K\"orper einen solchen, im Endlichen abbrechenden Algorithmus hat, so lassen sich dieselben Schl\"usse ziehen.

Dies findet statt bei dem K\"orper, der sich aus $R$ mittelst der dritten Einheitswurzel $\varrho = \dfrac{-1+\sqrt{-3}}{2}$ ableiten l\"asst und mit $R(\varrho)$ oder auch mit $R(\sqrt{-3})$ bezeichnet werden kann. Dieser K\"orper besteht aus dem Inbegriff aller Zahlen der Form $\xi = x+\varrho y$, worin $x$, $y$ rationale Zahlen sind, und die Norm einer solchen Zahl ist
\[
\begin{aligned}
N(\xi) &= (x+\varrho y)(x+\varrho^2 y) = x^2 - xy + y^2\\
&= \frac{2x-y+y\sqrt{-3}}{2}\cdot\frac{2x-y-y\sqrt{-3}}{2} = \frac{(2x-y)^2+3y^2}{4}.
\end{aligned}
\]
Eine Zahl $\alpha = a+b\varrho$ heisst eine ganze Zahl in $R(\varrho)$, wenn $a$, $b$ ganze rationale Zahlen sind. Die Einheiten in diesem K\"orper ergeben sich aus
\[
N(\alpha) = a^2 - ab + b^2 = 1,
\]
oder
\[
(2a-b)^2 + 3b^2 = 4,
\]
die nur erf\"ullt ist f\"ur $b=0$, $a=\pm 1$ oder $b=\pm 1$, $a=\pm 1$, so dass man also hier sechs Einheiten hat
% [source-original mathematical omission/ambiguity, verified 2026-08-29] Solving
% a^2-ab+b^2=1 gives (a,b)=(+-1,0),(0,+-1),(1,1),(-1,-1). The printed
% condition omits the a=0,b=+-1 cases and leaves the signs in its second pair ambiguous;
% the following six-unit list itself is correct and the printed prose is preserved.
\[
\pm 1,\quad \pm\varrho,\quad \pm\varrho^2.
\]
Ein System associirter Zahlen ist hier
\[
\begin{aligned}
&\pm(a+b\varrho),\quad \pm(a\varrho+b\varrho^2) = \pm[-b+(a-b)\varrho]\\
&\pm(a\varrho^2+b) = \pm(b-a-a\varrho).
\end{aligned}
\]
Ist $\xi$ eine beliebige gebrochene Zahl in $R(\varrho)$, so kann man die ganze Zahl $\mu$ so bestimmen, dass in $\xi-\mu = x+\varrho y$ die
%CONT>>
